\section{Runbook từng bước}
\label{sec:runbook}

\begin{keybox}
\textbf{Nguồn lệnh.} Các khối lệnh đánh dấu \textbf{[RUNBOOK]} được chép nguyên từ \code{code/RUNBOOK.md}
(chỉ bỏ hoặc dịch sang tiếng Anh không dấu các chú thích tiếng Việt, vì listings không in được chữ có dấu).
Khối đánh dấu \textbf{[DỰNG TỪ ARGPARSE]} do sổ tay ghép từ các cờ có thật trong \code{code/scripts/*.py} theo đúng quy tắc RUNBOOK.
Khối đánh dấu \textbf{[GỢI Ý]} là lệnh ngoài pipeline (tải dữ liệu, git\ldots), bạn nên kiểm tra lại.
Nếu RUNBOOK được cập nhật, \textbf{RUNBOOK thắng}.
\end{keybox}

% ===========================================================================
\subsection{Môi trường trên máy cá nhân (CPU)}
\label{sec:env-local}

Dùng để chạy kiểm thử và sửa code. Không cần GPU. Cần Python $\ge 3.9$ (\code{pyproject.toml}).

\textbf{[GỢI Ý]} (macOS/Linux; Windows dùng \code{.venv\textbackslash Scripts\textbackslash activate}):
\begin{lstlisting}[language=bash]
cd uit-cs315-lung-nodule-cxr-vlm/code
python3 -m venv .venv
source .venv/bin/activate
pip install -e ".[dev]"
pytest -q
\end{lstlisting}

\textbf{Kết quả mong đợi}: dòng cuối có \code{75 passed} (khi đã có torch, xem hộp dưới). Các nhóm phụ thuộc trong \code{pyproject.toml}:
\code{dev} = pytest; \code{torchcpu} = torch (CPU); \code{gpu} = torch, transformers, peft, accelerate, torchxrayvision, ultralytics.
Kiểm thử CPU không cần nhóm \code{gpu}.

\begin{warnbox}
\code{tests/test\_train.py} và một test trong \code{tests/test\_vlm.py} dùng \code{pytest.importorskip("torch")}: nếu máy chưa có \code{torch},
các test đó bị \emph{bỏ qua} (skipped) thay vì chạy. Để thấy đủ 75 passed, cài thêm torch bản CPU:
\code{pip install -e ".[dev,torchcpu]"} rồi chạy lại \code{pytest -q}.
\end{warnbox}

% ===========================================================================
\subsection{Clone repo vào Google Drive}
\label{sec:clone}

Colab đọc code và dữ liệu từ Google Drive. Repo là private, nên bạn cần (1) được mời làm collaborator, (2) quyền xác thực GitHub:
một \textbf{Personal Access Token} \en{PAT} hoặc \code{gh auth login} (RUNBOOK §0).

\textbf{[RUNBOOK] Lần đầu tiên} (chạy một lần duy nhất, trong một ô Colab): clone repo vào \code{/content/drive/MyDrive}.
\begin{lstlisting}[language=Python]
from google.colab import drive
drive.mount('/content/drive')
%cd /content/drive/MyDrive
!git clone https://github.com/minhdang900/uit-cs315-lung-nodule-cxr-vlm.git
\end{lstlisting}
Kết quả: bản clone nằm ở \code{/content/drive/MyDrive/uit-cs315-lung-nodule-cxr-vlm}. Mỗi phiên sau đó đặt
\code{REPO} bằng đường dẫn này (§\ref{sec:session}). Nếu bạn clone vào chỗ khác, chỉ cần đổi giá trị \code{REPO}.

\begin{warnbox}
\textbf{Xác thực khi clone repo private [GỢI Ý].} \code{git clone} sẽ hỏi tên người dùng và mật khẩu; mật khẩu là PAT, không phải mật khẩu GitHub.
Nếu ô Colab không cho nhập, đặt PAT vào Colab Secrets (ví dụ tên \code{GH\_TOKEN}) và clone bằng
\code{https://\{tok\}@github.com/...} với \code{tok = userdata.get("GH\_TOKEN")}. Sau đó chạy
\code{git remote set-url origin https://github.com/minhdang900/uit-cs315-lung-nodule-cxr-vlm.git} trong thư mục repo
để token không nằm lại trong \code{.git/config}. Không bao giờ dán token thẳng vào notebook.
\end{warnbox}

\textbf{Cách khác [GỢI Ý]}: trên máy có Google Drive for Desktop, clone vào thư mục \code{My Drive} đã đồng bộ; kết quả giống hệt.

\textbf{Thư mục dữ liệu.} Trong bản clone có \code{data/README.md}. Dữ liệu nằm ở \code{\$REPO/data/raw/} (tải về) và
\code{\$REPO/data/processed/} (do \code{prepare\_*} tạo). Cả hai bị gitignore; chỉ \code{data/README.md} vào git.

% ===========================================================================
\subsection{Quyền tải MedGemma và HF token}
\label{sec:hf}

\begin{enumerate}
\item Tạo tài khoản \url{https://huggingface.co}.
\item Mở trang \code{google/medgemma-4b-it}. Đọc và chấp nhận điều khoản \textbf{Health AI Developer Foundations (HAI-DEF)}. Chờ được cấp quyền (thường ngay lập tức, có thể lâu hơn).
\item Vào \emph{Settings $\rightarrow$ Access Tokens}, tạo token (quyền \emph{Read} là đủ để tải model).
\item Trong Colab: biểu tượng chìa khóa (\emph{Secrets}) ở thanh bên trái $\rightarrow$ \emph{Add new secret}, tên \code{HF\_TOKEN}, dán token, bật \emph{Notebook access}.
\end{enumerate}
Qwen2.5-VL-3B-Instruct không cần chấp nhận điều khoản riêng, nhưng vẫn tải qua cùng token.

\begin{dangerbox}
Không bao giờ dán token vào notebook, code, hay commit. Nếu lỡ lộ token: vào Hugging Face thu hồi \en{revoke} ngay và tạo token mới.
\end{dangerbox}

% ===========================================================================
\subsection{Colab và GPU}
\label{sec:colab}

\begin{itemize}
\item Cần \textbf{Colab Pro} (hoặc gói trả phí tương đương) để chọn GPU \textbf{L4} hoặc \textbf{A100}.
      \emph{Runtime $\rightarrow$ Change runtime type $\rightarrow$ GPU}.
\item \textbf{Không dùng T4.} MedGemma cần bf16 thật (compute capability $\ge 8$). Trên T4 (7.5), \code{load\_base} ném lỗi
      \code{RuntimeError: MedGemma needs a bf16-capable GPU (L4/A100); fp16 overflows}.
\item Colab tính phí theo ``compute unit''. Ngân sách thật phải đo bằng kiểm tra mục 7 và pilot E3, rồi ghi vào \code{outputs/budget.md}.
\end{itemize}

\begin{unsurebox}
Người bàn giao chưa chạy GPU lần nào, nên \textbf{chưa có số giờ GPU hay số compute unit thực tế}. Spec §7 chỉ đặt quy tắc:
nếu một lượt MedGemma vượt 3 giờ thì cắt giảm theo thứ tự (xem §\ref{sec:budget}).
\end{unsurebox}

% ===========================================================================
\subsection{Tải dữ liệu}
\label{sec:data}

\begin{dangerbox}
\textbf{Giấy phép dữ liệu.} Cả ba bộ chỉ được dùng cho nghiên cứu/học thuật theo điều khoản riêng của từng bộ.
\textbf{Không phân phối lại ảnh} (không commit, không đưa lên Drive công khai, không gửi cho người ngoài nhóm).
Phải trích dẫn nguồn trong báo cáo. Hãy đọc nguyên văn điều khoản trên trang của từng bộ; sổ tay này không thay thế điều khoản.
Spec ``Việc phải làm đầu tiên'' mục 1 yêu cầu xác nhận điều khoản trước khi dùng.
\end{dangerbox}

Bố cục dữ liệu (theo \code{data/README.md}): tải dữ liệu thô vào \code{data/raw/} (gợi ý mỗi bộ một thư mục con:
\code{data/raw/vinbig/}, \code{data/raw/nih/}, \code{data/raw/jsrt/}); \code{prepare\_*} ghi kết quả vào \code{data/processed/}.
Cả hai bị gitignore. Nếu bộ thô quá lớn so với dung lượng Drive (ví dụ VinBigData DICOM), có thể tải vào ổ tạm \code{/content/raw/}
của Colab (mất khi hết phiên) và chỉ để \code{data/processed/} trên Drive.

\subsubsection{VinBigData (huấn luyện, validation, test trong phân phối)}

\begin{itemize}
\item Nguồn: cuộc thi Kaggle \code{vinbigdata-chest-xray-abnormalities-detection}. Phải bấm chấp nhận luật cuộc thi \en{rules} trước khi tải.
\item Cần file \code{train.csv} (một dòng một box của một bác sĩ: \code{image\_id, class\_name, rad\_id, x\_min, y\_min, x\_max, y\_max};
      code chấp nhận cả \code{rad\_ID}) và ảnh train.
\item Ảnh gốc là DICOM, tên \code{<image\_id>.dicom}. Bộ gốc rất lớn (cỡ hàng trăm GB, kiểm tra trên trang Kaggle).
\item \textbf{Bản PNG đã thu nhỏ được hỗ trợ}: \code{prepare\_vinbig.py} có \code{--png-dir} và \code{--meta-csv}.
      Bản PNG phải đi kèm một CSV chứa \code{image\_id} và \textbf{kích thước gốc} của ảnh (để quy đổi tọa độ box trong \code{train.csv}).
      Mặc định cột chiều cao là \code{dim0}, chiều rộng là \code{dim1} (đổi bằng \code{--h-col}, \code{--w-col}).
      Ảnh PNG phải tên \code{<image\_id>.png}.
\end{itemize}

\begin{unsurebox}
Sổ tay chưa kiểm chứng tên bộ PNG cụ thể trên Kaggle. Các bộ cộng đồng kiểu ``vinbigdata chest xray resized png'' thường có
\code{train\_meta.csv} với cột \code{dim0}, \code{dim1}; hãy mở file kiểm tra: (1) \code{dim0}/\code{dim1} là kích thước \emph{gốc}, không phải kích thước PNG;
(2) \code{train.csv} đi kèm còn tọa độ gốc. Nếu sai, box sẽ lệch: \code{draw\_boxes.py} ở E0 sẽ cho thấy.
\end{unsurebox}

\textbf{Tùy chọn: tập test chính thức VinDr-CXR.} \code{prepare\_vinbig.py} có \code{--official-test-csv} và \code{--official-test-dicom}.
Khi dùng, VinBigData Kaggle được chia 85/15 (train/val) và tập test chính thức trở thành \code{test}. Phần này đọc DICOM bắt buộc.
Spec §6 ưu tiên dùng tập này nếu truy cập được (3.000 ảnh, 176 ảnh Nodule/Mass).

\begin{unsurebox}
Tập test VinDr-CXR nằm trên PhysioNet (cần tài khoản được xác thực). Sổ tay chưa kiểm chứng file nhãn test có cùng cột với \code{train.csv}
(đặc biệt cột \code{rad\_id}) hay không. Code gom box theo \code{rad\_id}; hãy kiểm tra trước khi chạy.
\end{unsurebox}

\subsubsection{NIH ChestX-ray14 (test ngoài phân phối 1)}

\begin{itemize}
\item Nguồn: NIH Clinical Center, \url{https://nihcc.app.box.com/v/ChestXray-NIHCC} (12 gói \code{images\_001..012.tar.gz}, tổng $\approx$45\,GB). Các link \code{nihcc.box.com/shared/static/...} cũ (trong script \code{batch\_download\_zips.py}) đã hỏng; tải theo \code{file\_id} như trong \code{scripts/download\_nih.sh}.
\item Cách nhanh nhất: \code{bash scripts/download\_nih.sh ../data/raw/nih}. Script tải 2 file CSV, tạo symlink \code{Data\_Entry\_2017.csv} $\rightarrow$ \code{Data\_Entry\_2017\_v2020.csv}, tính trước danh sách ảnh manifest cần (656 ảnh: 164 dương, 492 âm với \code{--neg-ratio 3}, \code{seed=0}), rồi đọc lần lượt từng gói và chỉ giữ các ảnh đó. Vẫn phải đọc đủ 45\,GB, nhưng chỉ lưu khoảng 0{,}3\,GB. Đổi \code{--neg-ratio} thì phải chạy lại script.
\item \code{prepare\_nih.py --raw-dir} cần trong thư mục đó (code tìm đệ quy):
  \begin{itemize}
  \item \code{Data\_Entry\_2017.csv} (cột dùng: \code{Image Index}, \code{Finding Labels}, \code{Patient ID});
  \item \code{BBox\_List\_2017.csv} (cột: \code{Image Index}, \code{Finding Label}, \code{Bbox [x}, \code{y}, \code{w}, \code{h]});
  \item ảnh \code{*.png} (1024$\times$1024), có thể nằm trong các thư mục con.
  \end{itemize}
\item \textbf{Phải tải toàn bộ ảnh}, không chỉ ảnh có box: ảnh âm được rút ngẫu nhiên (\code{seed=0}) từ \emph{mọi} ảnh ``No Finding'' của bệnh nhân không có ảnh dương,
      số lượng gấp \code{--neg-ratio} (mặc định 3) lần số ảnh dương. Thiếu ảnh sẽ gây \code{KeyError}.
\end{itemize}

\subsubsection{JSRT (test ngoài phân phối 2, có độ mờ)}

\begin{itemize}
\item Nguồn ảnh (chính thức, miễn phí): \url{http://imgcom.jsrt.or.jp/download/}, bản DICOM: \code{Nodule154images.zip} (806\,MB) và \code{NonNodule93images.zip} (484\,MB). Trang cũ \code{db.jsrt.or.jp} đã ngừng hoạt động.
\item File thông tin lâm sàng (tọa độ, kích thước, độ mờ) chính thức chỉ dành cho người đăng ký, mà trang đăng ký đã chết. Đồ án dùng bản sao \code{CLNDAT\_EN.txt} và \code{CNNDAT\_EN.TXT} trong repo HybridGNet (\url{https://github.com/ngaggion/HybridGNet}); chỉ dùng cho nghiên cứu, không commit (quyết định D28).
\item Cách nhanh nhất: \code{bash scripts/download\_jsrt.sh ../data/raw/jsrt} (tải ảnh, giải nén, tải 2 file lâm sàng).
\item 247 ảnh: 154 có nốt (\code{JPCLN*}), 93 không có nốt (\code{JPCNN*}); 2048$\times$2048, 0{,}175\,mm/pixel, MONOCHROME2.
\item \code{prepare\_jsrt.py} nhận cả bản \code{.IMG} cũ lẫn bản DICOM mới (ưu tiên \code{.IMG} nếu có), và nhận tên file lâm sàng \code{CLNDAT\_EN.txt}/\code{CLNDAT\_E.TXT}, \code{NNDAT\_EN.txt}/\code{CNNDAT\_EN.TXT}/\code{CNNDAT\_E.TXT} (không phân biệt hoa thường, tìm đệ quy).
\end{itemize}

\begin{keybox}
\textbf{Đã kiểm chứng (2026-11-13):} chạy \code{prepare\_jsrt.py} trên bản DICOM cho 154 ảnh dương, 93 ảnh âm; \code{draw\_boxes.py} cho thấy box nằm đúng nốt (ví dụ JPCLN001 thùy trên trái, JPCLN004 thùy trên phải), ảnh không bị đảo màu.
\end{keybox}

\subsubsection{Lệnh tải gợi ý}

\textbf{[GỢI Ý]} Kaggle CLI cần file API token \code{kaggle.json} (Kaggle $\rightarrow$ Settings $\rightarrow$ API). Không commit file này.
\begin{lstlisting}[language=bash]
pip install -q kaggle
mkdir -p ~/.kaggle && cp kaggle.json ~/.kaggle/ && chmod 600 ~/.kaggle/kaggle.json
kaggle competitions download -c vinbigdata-chest-xray-abnormalities-detection -p ../data/raw/vinbig
\end{lstlisting}

% ===========================================================================
\subsection{Chuẩn bị dữ liệu (phần đầu E0)}
\label{sec:prepare}

Chạy từ \code{code/}, với \code{M=../data/processed}. Các script này chỉ cần CPU (có thể chạy trên Colab CPU hoặc máy cá nhân có đủ dữ liệu).
RUNBOOK chỉ liệt kê cờ bắt buộc; dưới đây là lệnh đầy đủ \textbf{[DỰNG TỪ ARGPARSE]}, đường dẫn \code{../data/raw/...} theo bố cục ở trên (đổi nếu bạn để dữ liệu chỗ khác):

\begin{lstlisting}[language=bash]
%%bash
export M=../data/processed
# VinBigData, PNG variant (needs original sizes in --meta-csv)
python scripts/prepare_vinbig.py --train-csv ../data/raw/vinbig/train.csv --png-dir ../data/raw/vinbig/train --meta-csv ../data/raw/vinbig/train_meta.csv --out-root $M
# ... OR VinBigData, DICOM (sizes are read from the DICOM headers)
# python scripts/prepare_vinbig.py --train-csv ../data/raw/vinbig/train.csv --dicom-dir ../data/raw/vinbig/train --out-root $M
python scripts/prepare_nih.py --raw-dir ../data/raw/nih --out-root $M
python scripts/prepare_jsrt.py --raw-dir ../data/raw/jsrt --out-root $M
\end{lstlisting}

\textbf{Kết quả mong đợi}
\begin{itemize}
\item Trong \code{data/processed/}: \code{manifest\_vinbig.csv}, \code{manifest\_nih.csv}, \code{manifest\_jsrt.csv} và các thư mục \code{vinbig/}, \code{nih/}, \code{jsrt/} chứa PNG 896$\times$896 xám.
\item \code{prepare\_vinbig.py} in bảng \code{count}/\code{sum} theo split: tỷ lệ \code{train:val:test} xấp xỉ 70:10:20, tỷ lệ dương trong mỗi split gần bằng nhau (chia phân tầng).
\item \code{prepare\_nih.py} in số ảnh nhãn 1 và nhãn 0; nhãn 0 gấp 3 lần nhãn 1 (nếu đủ ảnh âm).
\item \code{prepare\_jsrt.py} in \code{1: 154}, \code{0: 93}.
\end{itemize}

\begin{unsurebox}
Spec ghi 410/15.000 ảnh Nodule/Mass (theo bài báo VinDr-CXR). Code coi một ảnh là dương nếu \emph{ít nhất một} bác sĩ khoanh Nodule/Mass,
và coi \emph{mọi} ảnh còn lại là âm (kể cả ảnh có bệnh khác, tức ``âm khó''). Vì vậy số ca dương thật có thể khác 410. Ghi con số thật vào E0 và báo cáo.
\end{unsurebox}

% ===========================================================================
\subsection{Mỗi phiên Colab (RUNBOOK §0)}
\label{sec:session}

\textbf{[RUNBOOK]} Ô Python đầu tiên của mỗi phiên. Đổi \code{REPO} nếu bạn clone vào chỗ khác (§\ref{sec:clone}):
\begin{lstlisting}[language=Python]
from google.colab import drive, userdata
drive.mount('/content/drive')
REPO = "/content/drive/MyDrive/uit-cs315-lung-nodule-cxr-vlm"
%cd {REPO}/code
import os; os.environ["HF_TOKEN"] = userdata.get("HF_TOKEN")   # from Colab Secrets
!pip install -q -e ".[gpu]"
!huggingface-cli login --token $HF_TOKEN
!mkdir -p /content/data && rsync -a ../data/processed/ /content/data/   # read images from local disk
\end{lstlisting}

\textbf{Giải thích từng dòng.} Gắn Drive; đặt \code{REPO} và chuyển vào \code{\$REPO/code}; nạp token vào biến môi trường; cài gói cùng nhóm \code{gpu};
đăng nhập Hugging Face; chép ảnh đã xử lý từ Drive sang ổ local \code{/content/data} (đọc từ Drive rất chậm).

Biến shell không giữ qua các dòng \code{!}. Vì vậy mọi ô chạy lệnh phải là ô \code{\%\%bash} và bắt đầu bằng \textbf{[RUNBOOK]}:
\begin{lstlisting}[language=bash]
%%bash
export M=../data/processed D=/content/data O=../outputs
\end{lstlisting}
\code{M}: manifest (trên Drive). \code{D}: gốc ảnh (ổ local). \code{O}: kết quả (trên Drive, gitignore).
Các khối lệnh bash bên dưới đều là nội dung đặt \emph{sau} dòng \code{export} này.

% ===========================================================================
\subsection{Bảy kiểm tra bắt buộc (phiên GPU đầu tiên)}
\label{sec:prechecks}

\begin{dangerbox}
Làm lần lượt, \textbf{trước} E1--E5. Mục nào không đạt thì \textbf{dừng và sửa}. Code GPU mới chỉ được kiểm tra cú pháp; đây là lần đầu nó chạy thật.
Ghi kết quả từng mục (đạt/không, con số đo được) vào \code{outputs/budget.md} hoặc một issue trên GitHub.
\end{dangerbox}

\begin{figure}[H]
\centering
\begin{tikzpicture}[node distance=0.25cm, every node/.style={font=\scriptsize}]
\node[obox,minimum width=3.2cm] (k1) at (0,0) {\textbf{1}\\Tokenizer yes/no};
\node[obox,minimum width=3.2cm] (k2) at (3.7,0) {\textbf{2}\\build\_example};
\node[obox,minimum width=3.2cm] (k3) at (7.4,0) {\textbf{3}\\Loss pilot 160 mẫu};
\node[obox,minimum width=3.2cm] (k4) at (11.1,0) {\textbf{4}\\Loại GPU};
\node[obox,minimum width=3.2cm] (k5) at (0,-1.7) {\textbf{5}\\Smoke YOLO};
\node[obox,minimum width=3.2cm] (k6) at (3.7,-1.7) {\textbf{6}\\Tốc độ mặt nạ phổi};
\node[obox,minimum width=3.2cm] (k7) at (7.4,-1.7) {\textbf{7}\\Tốc độ VLM};
\node[gbox,minimum width=3.2cm] (go) at (11.1,-1.7) {E1--E5};
\draw[arr] (k1)--(k2); \draw[arr] (k2)--(k3); \draw[arr] (k3)--(k4);
\draw[arr] (k4.south) -- ++(0,-0.3) -| (k5.north);
\draw[arr] (k5)--(k6); \draw[arr] (k6)--(k7); \draw[arr] (k7)--(go);
\end{tikzpicture}
\caption{Bảy kiểm tra trước khi chạy đủ. Gợi ý: làm mục 4 đầu tiên, vì nếu GPU sai thì mọi mục khác đều vô nghĩa.}
\label{fig:prechecks}
\end{figure}

\subsubsection*{Mục 1 -- Tokenizer cho điểm yes/no}
\textbf{[RUNBOOK]} Chạy cho \code{medgemma}, rồi lặp lại cho \code{qwen}:
\begin{lstlisting}[language=Python]
from nodulecxr.vlm.hf import load_base
from nodulecxr.vlm.chat import build_chat_text
from nodulecxr.prompts import QUESTION_LOC, ANSWER_PREFIX, build_target
_, processor = load_base("medgemma")   # repeat with "qwen"
tok = processor.tokenizer
base = build_chat_text(processor, QUESTION_LOC) + ANSWER_PREFIX
for w in ("yes", "no"):
    print(w, tok.convert_ids_to_tokens(tok(base + w, add_special_tokens=False)["input_ids"])[-8:])
tgt = build_chat_text(processor, QUESTION_LOC, build_target(1, [[1, 2, 3, 4]]))
print("target", tok.convert_ids_to_tokens(tok(tgt, add_special_tokens=False)["input_ids"]))
\end{lstlisting}
\textbf{Đạt khi}: (a) sau điểm phân kỳ giữa \code{yes} và \code{no}, token chỉ là chính từ trả lời, không gộp dấu nháy (\code{"yes} hay \code{yes"} là \textbf{sai});
(b) trong target huấn luyện, ngay sau \code{ANSWER\_PREFIX} là đúng token \code{yes}.
\textbf{Vì sao}: nếu dấu nháy bị gộp vào token, điểm $P(\text{yes})$ so sai hai chuỗi và mọi AUC đều sai.

\subsubsection*{Mục 2 -- Một lần \texttt{build\_example} trên ảnh thật}
\textbf{[RUNBOOK]} (dùng tiếp \code{processor}, \code{QUESTION\_LOC}, \code{build\_target} từ mục 1):
\begin{lstlisting}[language=Python]
from nodulecxr.records import read_manifest
from nodulecxr.images import load_image
from nodulecxr.train.sft_data import build_example
df = read_manifest("../data/processed/manifest_vinbig.csv"); r = df[df.label == 1].iloc[0]
ex = build_example(processor, load_image("/content/data", r.path), QUESTION_LOC, build_target(1, [[1, 2, 3, 4]]))
print({k: v.shape for k, v in ex.items()})
\end{lstlisting}
\textbf{Đạt khi}: in ra các shape, không có \code{ValueError: chat template: prompt is not a prefix of prompt+answer}.
Làm cho cả hai model.

\subsubsection*{Mục 3 -- Loss đầu tiên của pilot 160 mẫu}
RUNBOOK: thêm \code{--max-train 160 --epochs 1} vào một lệnh \code{train\_lora.py} của E3. Ví dụ \textbf{[DỰNG TỪ ARGPARSE]}
(thư mục \code{--out} tùy bạn, miễn không trùng \code{C1/C2/C3}):
\begin{lstlisting}[language=bash]
python scripts/train_lora.py --model medgemma --mode loc --manifest $M/manifest_vinbig.csv --data-root $D --seed 0 --max-train 160 --epochs 1 --out $O/lora/pilot
\end{lstlisting}
160 mẫu, grad-accum 8 $\rightarrow$ 20 bước tối ưu; \code{loss} được ghi ở bước 10 và 20 (\code{logging\_steps=10}).
\textbf{Đạt khi}: \code{loss} dòng đầu khoảng \textbf{0,5--3}. Nếu khoảng \textbf{10--25}: gradient accumulation bị chia hai lần --- dừng và sửa.
Đọc thêm \code{train\_loss} cuối và \code{timing.json} (giây huấn luyện) để ước tính thời gian một lượt đầy đủ.

\subsubsection*{Mục 4 -- Loại GPU}
\textbf{[RUNBOOK]}
\begin{lstlisting}[language=Python]
import torch; print(torch.cuda.get_device_name(), torch.cuda.get_device_capability())
\end{lstlisting}
\textbf{Đạt khi}: L4 hoặc A100, số đầu của capability $\ge 8$. Đừng tin \code{torch.cuda.is\_bf16\_supported()}: T4 (7.5) cũng trả \code{True} nhờ giả lập.

\subsubsection*{Mục 5 -- Smoke YOLO 1 epoch}
\textbf{[RUNBOOK]} Cách 1 (\code{<work>} là thư mục đã có \code{yolo\_data\_seed0/data.yaml}, ví dụ do cách 2 tạo):
\begin{lstlisting}[language=Python]
from ultralytics import YOLO
m = YOLO("yolov8s.pt"); m.train(data="<work>/yolo_data_seed0/data.yaml", imgsz=896, epochs=1, project="/content/yolo_smoke", name="s", exist_ok=True)
print(m.trainer.best)   # file must exist
\end{lstlisting}
\textbf{[RUNBOOK]} Cách 2, dùng đủ cờ bắt buộc của \code{train\_yolo.py}:
\begin{lstlisting}[language=bash]
%%bash
export M=../data/processed D=/content/data O=../outputs
python scripts/train_yolo.py --manifests $M/manifest_vinbig.csv $M/manifest_nih.csv $M/manifest_jsrt.csv --data-root $D --epochs 1 --work /content/yolo_smoke --out /content/smoke.jsonl
ls /content/yolo_smoke/seed0/weights/best.pt
\end{lstlisting}
Cách 2 sau khi train còn dự đoán toàn bộ val, test, test\_nih, test\_jsrt nên chậm; có thể dừng ô khi \code{best.pt} đã xuất hiện.

\subsubsection*{Mục 6 -- Thời gian mặt nạ phổi}
\textbf{[RUNBOOK]} (PSPNet của torchxrayvision chạy trên CPU; nhãn in đã đổi sang ASCII):
\begin{lstlisting}[language=Python]
import sys, time; sys.path.insert(0, "scripts")
from run_audit import lung_mask_xrv
from nodulecxr.records import read_manifest
from nodulecxr.images import load_image
df = read_manifest("../data/processed/manifest_jsrt.csv").head(10)
lung_mask_xrv(load_image("/content/data", df.iloc[0].path))   # warm-up, not timed
t = time.time()
for r in df.itertuples(): lung_mask_xrv(load_image("/content/data", r.path))
print((time.time() - t) / len(df), "s/img")
\end{lstlisting}
Nếu quá chậm: dùng \code{--lung-mask heuristic} ở E5 và ghi quyết định vào \code{\$O/budget.md}.

\subsubsection*{Mục 7 -- Thời gian suy luận mỗi ảnh của VLM}
\textbf{[RUNBOOK]} 50 ảnh val mỗi model, config \code{pilot}, ghi ra ngoài \code{\$O/preds/}:
\begin{lstlisting}[language=bash]
for m in medgemma qwen; do
  python scripts/predict_vlm.py --model $m --manifests $M/manifest_vinbig.csv --splits val --limit 50 --data-root $D --config pilot --max-new-tokens 256 --out /content/pilot_$m.jsonl
done
\end{lstlisting}
Đọc \code{s/it} ở thanh tqdm (mỗi \code{it} = một ảnh, gồm cả tính điểm và sinh văn bản). Ghi vào \code{\$O/budget.md}.
Chi phí C0 $\approx$ giây/ảnh $\times$ số ảnh của \code{val test test\_nih test\_jsrt} $\times$ 3 run $\times$ 2 model.
Nhân tiện, mở \code{/content/pilot\_medgemma.jsonl} xem \code{raw\_text} và \code{parse\_ok}: MedGemma có xuất JSON và box không? (spec ``Việc phải làm đầu tiên'' mục 3).

% ===========================================================================
\subsection{Chạy các thí nghiệm}
\label{sec:exp-run}

\subsubsection{E0 -- Dữ liệu}
Sau 3 script \code{prepare\_*} (§\ref{sec:prepare}), \textbf{[RUNBOOK]}:
\begin{lstlisting}[language=bash]
for d in vinbig nih jsrt; do python scripts/draw_boxes.py --manifest $M/manifest_$d.csv --data-root $M --out $O/qa/$d -n 8; done
python scripts/data_stats.py --manifests $M/manifest_vinbig.csv $M/manifest_nih.csv $M/manifest_jsrt.csv --out $O/e0_data.md
\end{lstlisting}
\textbf{Kiểm tra}: mở \code{outputs/qa/\{vinbig,nih,jsrt\}/} bằng mắt; box đỏ phải nằm đúng trên nốt. Đọc \code{outputs/e0\_data.md}.
\textbf{Đầu ra}: \code{\$O/e0\_data.md}, \code{\$O/e0\_data\_edges.json} (ngưỡng tertile, dùng ở báo cáo và E5).

\subsubsection{E1 -- Baseline CNN (3 seed mỗi loại)}
\textbf{[RUNBOOK]}
\begin{lstlisting}[language=bash]
for s in 0 1 2; do
  python scripts/train_densenet.py --manifests $M/manifest_vinbig.csv $M/manifest_nih.csv $M/manifest_jsrt.csv --data-root $D --seed $s --out $O/preds/B1.jsonl
  python scripts/train_yolo.py --manifests $M/manifest_vinbig.csv $M/manifest_nih.csv $M/manifest_jsrt.csv --data-root $D --seed $s --work /content/yolo --out $O/preds/B2.jsonl
  mkdir -p $O/yolo/seed$s/weights && cp /content/yolo/seed$s/weights/best.pt $O/yolo/seed$s/weights/best.pt
done
\end{lstlisting}
\textbf{Kiểm tra}: log DenseNet in \code{epoch k val AUC ...} mỗi epoch; \code{B1.jsonl}, \code{B2.jsonl} có dòng cho mọi ảnh không thuộc train;
\code{\$O/yolo/seed\{0,1,2\}/weights/best.pt} tồn tại trên Drive. Dòng \code{cp} bắt buộc: \code{/content} mất khi hết phiên, và E5 cần \code{best.pt}.

\subsubsection{E2 -- Zero-shot C0}
\textbf{[RUNBOOK]}
\begin{lstlisting}[language=bash]
for m in medgemma qwen; do
  python scripts/predict_vlm.py --model $m --manifests $M/manifest_vinbig.csv $M/manifest_nih.csv $M/manifest_jsrt.csv --splits val test test_nih test_jsrt --data-root $D --config C0-$m --run 0 --temperature 0 --max-new-tokens 256 --out $O/preds/C0-$m.jsonl
  for r in 1 2; do
    python scripts/predict_vlm.py --model $m --manifests $M/manifest_vinbig.csv $M/manifest_nih.csv $M/manifest_jsrt.csv --splits val test test_nih test_jsrt --data-root $D --config C0-$m --run $r --temperature 0.7 --max-new-tokens 256 --out $O/preds/C0-$m.jsonl
  done
done
\end{lstlisting}
Script tự bỏ qua ảnh đã có trong file (cùng \code{config, seed, run}); Colab ngắt thì chạy lại \emph{đúng lệnh cũ}.
Bảng chính chỉ dùng run 0 (nhiệt độ 0); run 1--2 dùng cho \code{run\_agreement.csv}.

\subsubsection{E3 -- LoRA MedGemma C1--C3 (3 seed)}
Trước tiên: pilot ở kiểm tra mục 3. Sau đó \textbf{[RUNBOOK]} huấn luyện:
\begin{lstlisting}[language=bash]
for s in 0 1 2; do
  python scripts/train_lora.py --model medgemma --mode cls --manifest $M/manifest_vinbig.csv --data-root $D --seed $s --out $O/lora/C1-s$s
  python scripts/train_lora.py --model medgemma --mode loc --manifest $M/manifest_vinbig.csv --data-root $D --seed $s --out $O/lora/C2-s$s
  python scripts/train_lora.py --model medgemma --mode loc --balance --pos-weight 2.0 --manifest $M/manifest_vinbig.csv --data-root $D --seed $s --out $O/lora/C3-s$s
done
\end{lstlisting}

Dự đoán. \textbf{[RUNBOOK]} ví dụ cho C3:
\begin{lstlisting}[language=bash]
for s in 0 1 2; do
  python scripts/predict_vlm.py --model medgemma --adapter $O/lora/C3-s$s/final --prompt loc --generate all --manifests $M/manifest_vinbig.csv $M/manifest_nih.csv $M/manifest_jsrt.csv --splits val test test_nih test_jsrt --data-root $D --config C3 --seed $s --out $O/preds/C3.jsonl
done
\end{lstlisting}
C1 và C2 \textbf{[DỰNG TỪ ARGPARSE]} theo quy tắc RUNBOOK (C1: \code{--prompt cls --generate all --max-new-tokens 16}; C2: \code{--prompt loc --generate all}):
\begin{lstlisting}[language=bash]
for s in 0 1 2; do
  python scripts/predict_vlm.py --model medgemma --adapter $O/lora/C1-s$s/final --prompt cls --generate all --max-new-tokens 16 --manifests $M/manifest_vinbig.csv $M/manifest_nih.csv $M/manifest_jsrt.csv --splits val test test_nih test_jsrt --data-root $D --config C1 --seed $s --out $O/preds/C1.jsonl
  python scripts/predict_vlm.py --model medgemma --adapter $O/lora/C2-s$s/final --prompt loc --generate all --manifests $M/manifest_vinbig.csv $M/manifest_nih.csv $M/manifest_jsrt.csv --splits val test test_nih test_jsrt --data-root $D --config C2 --seed $s --out $O/preds/C2.jsonl
done
\end{lstlisting}
Với C1, \code{--generate none} chỉ cho điểm và để trống các cột \code{raw\_sens}, \code{no\_on\_pos}, \code{yes\_on\_neg}, \code{unparsed}; vì vậy nên dùng \code{--generate all}.

\textbf{Đường học} (RUNBOOK): chạy \code{predict\_vlm.py --splits val --generate none} cho từng \code{checkpoint-*} của seed 0, với
\code{--config <cfg>-ep<k>} và cùng \code{--prompt} của cấu hình nguồn.

\begin{warnbox}
Nếu ghi các dự đoán đường học vào \code{\$O/preds/}, \code{make\_report.py} sẽ đưa các dòng \code{C3-ep1}, \code{C3-ep2}\ldots{} vào \code{val\_selection.csv}
(chúng chỉ có split val nên không vào bảng test). Khi chọn cấu hình tốt nhất cho E4/E5, \textbf{chỉ so C1, C2, C3}; bỏ qua các dòng \code{-ep}.
\end{warnbox}

\subsubsection{Báo cáo trung gian (cổng chọn cấu hình)}
\textbf{[RUNBOOK]}
\begin{lstlisting}[language=bash]
python scripts/make_report.py --manifests $M/manifest_vinbig.csv $M/manifest_nih.csv $M/manifest_jsrt.csv --preds "$O/preds/*.jsonl" --edges-json $O/e0_data_edges.json --out $O/report
\end{lstlisting}
Mở \code{\$O/report/val\_selection.csv}. Tính trung bình \code{val\_sens90} qua seed cho C1, C2, C3. Cấu hình cao nhất là ``LoRA tốt nhất''.
\code{val\_auc} chỉ để tham khảo. Ghi lựa chọn và con số vào \code{docs/decisions.md}.

\textbf{[GỢI Ý]} đọc nhanh bằng pandas:
\begin{lstlisting}[language=Python]
import pandas as pd
v = pd.read_csv("../outputs/report/val_selection.csv")
print(v[v.config.isin(["C1", "C2", "C3"])].groupby("config").val_sens90.agg(["mean", "std", "count"]))
\end{lstlisting}

\subsubsection{E4 -- C4 trên Qwen2.5-VL-3B (1--3 seed)}
RUNBOOK: chạy lại \code{train\_lora.py --model qwen} với \emph{đúng các cờ} của cấu hình tốt nhất, rồi dự đoán với \code{--model qwen --config C4}
và đúng \code{--prompt}/\code{--generate} (cùng \code{--max-new-tokens}) của cấu hình nguồn.
Ví dụ \textbf{[DỰNG TỪ ARGPARSE]} nếu cấu hình tốt nhất là C3, seed 0:
\begin{lstlisting}[language=bash]
python scripts/train_lora.py --model qwen --mode loc --balance --pos-weight 2.0 --manifest $M/manifest_vinbig.csv --data-root $D --seed 0 --out $O/lora/C4-s0
python scripts/predict_vlm.py --model qwen --adapter $O/lora/C4-s0/final --prompt loc --generate all --manifests $M/manifest_vinbig.csv $M/manifest_nih.csv $M/manifest_jsrt.csv --splits val test test_nih test_jsrt --data-root $D --config C4 --seed 0 --out $O/preds/C4.jsonl
\end{lstlisting}
Nếu tốt nhất là C1: thay \code{--mode loc --balance --pos-weight 2.0} bằng \code{--mode cls}, và dự đoán với \code{--prompt cls --generate all --max-new-tokens 16}.

\subsubsection{E5 -- Kiểm toán che vùng}
Cấu hình: C0-medgemma, LoRA tốt nhất (adapter seed 0), B2 (seed 0). Hai tập:
\begin{itemize}
\item \code{test}, dùng \code{manifest\_vinbig.csv};
\item \code{test\_jsrt}, dùng \code{manifest\_jsrt.csv}.
\end{itemize}
Ngưỡng lấy từ \code{\$O/report/thresholds.json}, khóa \code{config|seed|run}. \textbf{[GỢI Ý]} đọc ngưỡng:
\begin{lstlisting}[language=bash]
python -c "import json; t=json.load(open('../outputs/report/thresholds.json')); print(t['C3|0|0']['t90'], t['C0-medgemma|0|0']['t90'], t['B2|0|0']['t90'])"
\end{lstlisting}
\textbf{[RUNBOOK]} (ví dụ LoRA tốt nhất là C3; thay \code{<t90...>} bằng số vừa đọc):
\begin{lstlisting}[language=bash]
# best LoRA (example: C3, seed 0; if C1, add --prompt cls)
python scripts/run_audit.py --scorer vlm --model medgemma --adapter $O/lora/C3-s0/final --manifest $M/manifest_jsrt.csv --split test_jsrt --data-root $D --threshold <t90> --edges-json $O/e0_data_edges.json --out $O/audit/C3_jsrt.json
# C0 zero-shot (no --adapter)
python scripts/run_audit.py --scorer vlm --model medgemma --manifest $M/manifest_jsrt.csv --split test_jsrt --data-root $D --threshold <t90 of C0-medgemma|0|0> --edges-json $O/e0_data_edges.json --out $O/audit/C0_jsrt.json
# B2 YOLO
python scripts/run_audit.py --scorer yolo --yolo-weights $O/yolo/seed0/weights/best.pt --manifest $M/manifest_jsrt.csv --split test_jsrt --data-root $D --threshold <t90 of B2|0|0> --edges-json $O/e0_data_edges.json --out $O/audit/B2_jsrt.json
\end{lstlisting}
Lặp lại cho VinBigData test: đổi thành \code{--manifest \$M/manifest\_vinbig.csv} và \code{--split test}, đổi tên \code{--out} (ví dụ \code{C3\_test.json}).
Thêm \code{--lung-mask heuristic} nếu mặt nạ PSPNet quá chậm (mục 6).

\textbf{Đọc kết quả}: file \code{.json} có \code{n}, \code{n\_orig\_yes}, \code{n\_t2\_ok}, \code{flip\_t1\_all}, \code{flip\_t1}, \code{flip\_t2}, \code{delta}, \code{yes\_rate\_t3},
cùng \code{by\_size} và (JSRT) \code{by\_subtlety}. \textbf{Số cần báo cáo là \code{delta}}, kèm \code{n\_t2\_ok}. File \code{.csv} cùng tên chứa điểm từng ảnh.

\subsubsection{Báo cáo cuối}
Chạy lại \code{make\_report.py} (cùng lệnh ở trên) sau khi có C4. Thư mục \code{\$O/report/} khi đó có đủ 8 file.

% ===========================================================================
\subsection{Ngân sách và cắt giảm}
\label{sec:budget}

Spec §7: lượt LoRA đầu tiên dùng để đo thời gian. \textbf{Nếu một lượt vượt 3 giờ}, cắt theo đúng thứ tự (RUNBOOK):

\begin{figure}[H]
\centering
\begin{tikzpicture}[node distance=0.35cm, every node/.style={font=\scriptsize}]
\node[obox,minimum width=3.0cm] (a) {1. \code{--neg-per-pos} 4 $\rightarrow$ 2\\(giữ mọi ảnh dương)};
\node[obox,right=of a,minimum width=3.2cm] (b) {2. Giảm độ phân giải\\chỉ Qwen: \code{min\_pixels}/\code{max\_pixels}\\trong \code{vlm/hf.py}};
\node[obox,right=of b,minimum width=2.4cm] (c) {3. \code{--epochs 2}};
\node[obox,right=of c,minimum width=2.4cm] (d) {4. C1 chỉ 1 seed};
\draw[arr] (a)--(b); \draw[arr] (b)--(c); \draw[arr] (c)--(d);
\end{tikzpicture}
\caption{Thứ tự cắt giảm. Bước 2 không áp dụng cho MedGemma (processor cố định 896 px) và không có cờ dòng lệnh.}
\label{fig:budget}
\end{figure}

Ghi \emph{mọi} quyết định vào \code{outputs/budget.md}: ngày, đo được bao nhiêu, cắt bước nào, áp cho cấu hình nào.
\code{outputs/} không vào git, nên tóm tắt các quyết định cắt giảm cũng phải chép sang \code{docs/decisions.md}.

% ===========================================================================
\subsection{Lỗi thường gặp}
\label{sec:errors}

\begin{center}
\small
\begin{longtable}{@{}P{5.0cm}P{9.6cm}@{}}
\toprule
\textbf{Triệu chứng} & \textbf{Nguyên nhân và cách sửa} \\
\midrule\endhead
\code{RuntimeError: MedGemma needs a bf16-capable GPU (L4/A100)} & Đang ở T4 hoặc CPU. Đổi runtime sang L4/A100. Kiểm tra lại bằng mục 4. \\
\code{401}/\code{403}, \code{GatedRepoError}, ``access to model is restricted'' & Chưa được cấp quyền MedGemma, hoặc token sai/thiếu. Kiểm tra: đã chấp nhận HAI-DEF trên đúng tài khoản; secret tên đúng \code{HF\_TOKEN}; đã bật \emph{Notebook access}; đã chạy ô §0. \\
\code{SecretNotFoundError} khi \code{userdata.get} & Chưa tạo secret, hoặc gõ sai tên. \\
Colab ngắt kết nối giữa chừng & \code{predict\_vlm.py}: chạy lại đúng lệnh, nó chạy tiếp từ ảnh chưa làm. \code{train\_lora.py}, \code{train\_densenet.py}, \code{train\_yolo.py}: \textbf{không} tự chạy tiếp; chạy lại sẽ huấn luyện từ đầu. B1/B2 sẽ ghi thêm dòng trùng vào JSONL; \code{make\_report.py} giữ dòng cuối nên vẫn đúng. Nếu một dòng JSONL cuối bị cắt dở, xóa dòng đó. \\
\code{CUDA out of memory} & Batch đã là 1 và đã bật gradient checkpointing. Thử: khởi động lại runtime (giải phóng model đã tải ở các ô kiểm tra), dùng A100 thay L4. Với Qwen có thể giảm \code{max\_pixels} trong \code{vlm/hf.py} (bước 2 của quy tắc ngân sách). Ghi mọi thay đổi vào \code{budget.md}. \\
\code{loss} ban đầu khoảng 10--25 & Gradient accumulation bị chia hai lần (mục 3). Dừng, không chạy E3; báo lại trong issue. \\
\code{ValueError: chat template: prompt is not a prefix\ldots} & Chat template của model/phiên bản transformers khác giả định. Dừng; ghi phiên bản \code{transformers}; không tự ý bỏ kiểm tra. \\
\code{KeyError} trong \code{prepare\_nih.py} & Thiếu ảnh NIH: phải tải đủ các gói ảnh. \\
\code{StopIteration} trong \code{prepare\_jsrt.py} & Thiếu file \code{.IMG} của một ảnh có trong file TXT. \\
\code{FileNotFoundError: ...CLNDAT\_EN.txt} & Sai tên hoặc vị trí file thông tin JSRT; phải nằm ngay trong \code{--raw-dir}. \\
Box vẽ ở E0 bị lệch & Kích thước gốc sai (\code{--meta-csv}, \code{--h-col}, \code{--w-col}). Sửa rồi xóa \code{data/processed/vinbig/} và chạy lại (script bỏ qua PNG đã có). \\
Chạy rất chậm khi đọc ảnh & Đang đọc từ Drive. Dùng \code{--data-root \$D} (\code{/content/data}) sau lệnh \code{rsync} ở §0. \\
\code{t90} là \code{Infinity} & Mô hình cho điểm gần như bằng nhau trên mọi ảnh val (chế độ báo bình thường cực đoan). Không sửa; báo cáo trung thực, sensitivity sẽ là 0. \\
Mất \code{best.pt} của YOLO & Quên dòng \code{cp} ở E1. Phải huấn luyện lại seed đó. \\
\bottomrule
\end{longtable}
\end{center}

% ===========================================================================
\subsection{Điền báo cáo từ kết quả}
\label{sec:fill-report}

Báo cáo nằm ở \code{report/}. Các ô \textcolor{badred}{[--]} màu đỏ chờ số. Biên dịch: \code{latexmk -xelatex main.tex} trong \code{report/}.

\begin{center}
\small
\begin{longtable}{@{}P{4.2cm}P{4.8cm}P{5.6cm}@{}}
\caption{Ánh xạ file kết quả $\rightarrow$ vị trí trong báo cáo.}\label{tab:fillmap}\\
\toprule
\textbf{Vị trí trong báo cáo} & \textbf{File nguồn} & \textbf{Cột / khóa} \\
\midrule\endfirsthead
\toprule
\textbf{Vị trí trong báo cáo} & \textbf{File nguồn} & \textbf{Cột / khóa} \\
\midrule\endhead
Bảng \code{tab:data-stats} (\code{04a\_du\_lieu.tex}) & \code{outputs/e0\_data.md} & \code{n}, \code{n\_pos}, \code{prevalence} theo split; tầng kích thước \\
Bảng \code{tab:main} (\code{04b\_ket\_qua.tex}) & \code{outputs/report/summary.csv} & \code{sens\_mean} $\pm$ \code{sens\_sd} (sens ở $t_{90}$), \code{auc\_roc\_mean}; theo \code{split} = \code{test}, \code{test\_nih}, \code{test\_jsrt} \\
CI 95\% kèm bảng chính & \code{summary\_per\_seed.csv} & \code{sens\_lo}, \code{sens\_hi}, \code{auc\_roc\_lo}, \code{auc\_roc\_hi} (theo seed) \\
Bảng \code{tab:loc} & \code{summary.csv} & \code{mean\_best\_iou\_mean}, \code{hit\_iou30\_mean}, \code{hit\_iou50\_mean} (\code{test}); \code{jsrt\_median\_dist\_mm\_mean} (\code{test\_jsrt}). Trừ \code{hit\_iou30}, các cột này không có trong \code{summary.md}. \\
Mục ``Chế độ báo bình thường'' & \code{summary.csv} & \code{no\_on\_pos}, \code{yes\_on\_neg}, \code{unparsed}, \code{raw\_sens}, \code{ece}, \code{balanced\_brier}; đường học từ dự đoán \code{-ep<k>} \\
Độ ổn định C0 (E2) & \code{run\_agreement.csv} & \code{agreement}, \code{n\_images} \\
Bảng \code{tab:audit-res} & \code{outputs/audit/*.json} & \code{flip\_t1}, \code{flip\_t2}, \code{delta} (\code{n\_t2\_ok}); cột T3 ``ảnh âm $\rightarrow$ no'' $= 1 -$ \code{yes\_rate\_t3} \\
Biểu đồ sens theo tầng & \code{strata.csv}; \code{audit/*.json} & \code{kind} (\code{size}/\code{subtlety}), \code{stratum}, \code{sensitivity}, \code{count}; \code{by\_size}, \code{by\_subtlety} \\
p-value so sánh cặp & \code{comparisons.csv} & \code{p\_sens\_holm}, \code{p\_spec\_holm}, \code{p\_auc\_holm} (seed 0, run 0) \\
Chọn cấu hình (phương pháp) & \code{val\_selection.csv} & \code{val\_sens90} trung bình qua seed \\
FROC, AUC-PR, PPV & \code{summary.csv} & \code{froc\_0.5fp\_mean}, \code{froc\_1fp\_mean}, \code{auc\_pr\_mean}, \code{ppv\_prev\_1pct\_mean} \\
Tóm tắt, Kết luận (\code{00\_tom\_tat.tex}, \code{05\_ket\_luan.tex}) & tổng hợp các bảng trên & trả lời Q1, Q2 \\
\bottomrule
\end{longtable}
\end{center}

\begin{keybox}
\textbf{Cách điền.} Mỗi lần điền số, ghi rõ trong commit: số lấy từ file nào, chạy ngày nào (ví dụ \code{docs: fill tab:main from outputs/report/summary.csv (2026-12-01)}).
Báo cáo số làm tròn 3 chữ số thập phân, dùng dấu phẩy thập phân theo văn phong tiếng Việt của báo cáo. Ô nào không có số (ví dụ C1 không có box) ghi ``---'' và giải thích.
\end{keybox}
